\subsection*{Degree-six recognition for every once-punctured elliptic curve}

\paragraph{Theorem.}
Let $E$ be an elliptic curve over a number field $K$, put
$U=E\setminus\{O\}$, and let $\ell$ be any prime. Set
\[
 P=\pi_1(U_{\overline K})^{(\ell)},\qquad
 \Pi=\pi_1(U)/\ker\bigl(\pi_1(U_{\overline K})\to P\bigr),
 \qquad p:\Pi\to G_K.
\]
Every nontrivial closed procyclic $A\subset P$ whose arithmetic
normalizer $N=N_\Pi(A)$ has open image under $p$ is contained in an
actual cusp inertia subgroup of $P$.

\paragraph{1. The exact primary input.}
Nakamura1996, Lemma4.5 and Corollary4.6,p.116, give the surjection
from the graded arithmetic Galois image of a fixed affine hyperbolic
curve to its universal arithmetic quotient. Equations(4.7.1)--(4.7.2),
p.117, identify this quotient and its action: in degree$m$ it is acted
on by the similitude character to the power$m/2$. For $(g,n)=(1,1)$,
degree6 is nonzero by Corollary4.2 and the degree6 entry of(4.3),p.115.
These are unconditional statements in the original article; the
standing prime in Section2,p.101 is any rational prime. The
$\operatorname{End}(J)=\mathbf Z$ condition in Theorem B,p.100,
concerns its additional centralizer-finiteness conclusion and is not
a premise of Lemma4.5 or Corollary4.6.

In explicit terms, put
\[
 G_K(m)=\{\sigma\in G_K:\rho_U(\sigma)\in\Gamma_{1,1}(m)\}.
\]
The degree-six image $\mathfrak a_6\subset
\operatorname{gr}_6\Gamma_{1,1}\otimes\mathbf Q_\ell$ has a nonzero
$G_K$-equivariant quotient
\[
 q_6:\mathfrak a_6\twoheadrightarrow Q_6,
 \qquad G_K\text{ acts on }Q_6\text{ by }\chi_\ell^3.
\]
Here $\mathfrak a_6$ is the rational span of the image of the actual
group $G_K(6)$ in $\Gamma_{1,1}(6)/\Gamma_{1,1}(7)$.
This is the quotient used below; no general normalizer assertion is
attributed to Theorem B.

\paragraph{2. The first Magnus degree.}
The finite-quotient detector of F.3.15,Steps1--2 applies to every
once-punctured elliptic curve. For clarity its precise mechanism is
as follows. A generator $u$ of $A\simeq\mathbf Z_\ell$ survives a
finite $\ell$-group quotient. The inverse image of its cyclic image
is an open subgroup whose torsion-free abelianization detects $u$.
It is the maximal pro-$\ell$ group of an actual finite etale cover;
normal closures of the relevant iterated $\ell$-covers prove that
identification. Descending the cover and intersecting its arithmetic
subgroup with $N$ gives the line of $u$ in actual curve homology,
without choosing a section of $N$. Its Frobenius weight is one or
two, by the proper-homology and boundary exact sequence.

The injective Magnus embedding and its rational logarithm give a
first nonzero term $u_d$ in the degree$d$ free Lie algebra on
$V=V_\ell(E)$. Geometric conjugation is trivial on that term, so the
normalizer exponent $\theta$ descends to $p(N)$ and satisfies
$\sigma(u_d)=\theta(\sigma)u_d$. The degree$d$ term has Frobenius
weight$d$. Comparing the same character with the detector shows
$d=1$ or2. These weight statements use arithmetic Frobenius on
homology, so $V_\ell(E)$ has eigenvalue absolute value$q^{1/2}$.

\paragraph{3. The only degree-one case needing an obstruction.}
Extend $K$ finitely so that all geometric endomorphisms of $E$ are
defined. Openness of the arithmetic normalizer image is preserved.
Write $H=p(N)$ after this replacement. Faltings' semisimplicity and
endomorphism theorem, original \S5,S\"atze3--4,p.360, give
\[
 \operatorname{End}_{G_L}(V_\ell E)
 =\operatorname{End}^0_L(E)\otimes\mathbf Q_\ell
\]
for every further finite extension $L/K$.
For a non-CM curve this algebra is $\mathbf Q_\ell$, which has no
rank-one projector on the two-dimensional $V$. For a CM curve at a
prime not split in its CM field, the algebra is a quadratic field,
which also has no nontrivial idempotent. Semisimplicity therefore
excludes a one-dimensional subrepresentation of any open arithmetic
subgroup in those cases.

The remaining case is split CM, where
\[
 V=\mathbf Q_\ell x\oplus\mathbf Q_\ell y,
 \qquad \rho_V=\chi_1\oplus\chi_2,
 \qquad \chi_1\chi_2=\chi_\ell.
\]
These are the only stable lines. The ratio $\chi_1/\chi_2$ has
infinite image on $H$. Otherwise a finite extension would make the
two characters equal, and the displayed endomorphism algebra would
become $M_2(\mathbf Q_\ell)$ instead of the split quadratic algebra,
contradicting Faltings' theorem and the full geometric CM
endomorphism algebra. Since $\mathbf Q_\ell$ has only finitely many
roots of unity, choose an actual $\tau\in H$ for which
$\chi_1(\tau)/\chi_2(\tau)$ is not a root of unity.

\paragraph{4. The character projection produces an actual arithmetic element.}
Use Nakamura1995,(4.3)--(4.4),p.214, to lift the degree-at-least-three
outer action canonically to a boundary-fixing representative.
The isomorphism of graded pieces there shows that an element of
$G_K(6)$ has representative in $\Gamma^*_{1,1}(6)$.
Its leading derivation $D$ has degree6 and is special:
\[
 D([x,y])=0.
\]
Let $\mathcal D_6$ denote this finite-dimensional space over
$\mathbf Q_\ell$; the integral action gives it a fixed
$\mathbf Z_\ell$-lattice containing the actual image. Set
\[
 M=\operatorname{im}\bigl(H\cap G_K(6)\longrightarrow
       \operatorname{gr}_6\Gamma^*_{1,1}\bigr).
\]
This is a compact, hence closed, additive subgroup of a finite free
$\mathbf Z_\ell$-module. It is a $\mathbf Z_\ell$-module: multiplication
by an $\ell$-adic integer is the limit of integer multiples, and
closedness keeps that limit in $M$. It is invariant under
$T=\operatorname{Ad}(\tau)$, since $G_K(6)$ is normal and
$\tau\in H$. Moreover $M\otimes\mathbf Q_\ell=\mathfrak a_6$;
restriction to the finite-index subgroup $H\cap G_K(6)$ does not
change the rational image.

The eigencharacters of $T$ on $\mathcal D_6\subset
\operatorname{Hom}(V,L_7(V))$ have the form
$\chi_1^a\chi_2^b$ with $a+b=6$. Because the chosen ratio is not
a root of unity, the eigenvalue
\[
 c_0=\chi_\ell(\tau)^3
\]
occurs precisely for the middle weight $(a,b)=(3,3)$.
For the finite set of distinct eigenvalues $c$ of $T$, the polynomial
\[
 e(T)=\prod_{c\ne c_0}\frac{T-c}{c_0-c}
\]
is the projection to this middle weight. Its coefficients belong
to $\mathbf Q_\ell$. Choose an integer $r\ge0$ such that all
coefficients of $\ell^r e(T)$ are in $\mathbf Z_\ell$.

By Step1 choose $m\in M$ with $q_6(m)\ne0$. Equivariance gives
$q_6(e(T)m)=q_6(m)$, since $T$ acts on the quotient by $c_0$.
Then
\[
 m_0=\ell^r e(T)m\in M,\qquad q_6(m_0)\ne0.
\]
Membership in $M$ follows from its $\mathbf Z_\ell$-module property
and invariance under all powers of $T$. By the definition of $M$,
there is an actual
\[
 \sigma\in H\cap G_K(6)
\]
whose degree-six derivation is $m_0$. Thus its first nonzero action
is the nonzero special derivation $D=m_0$ in middle weight $(3,3)$.
The use of closed image and cleared denominators is essential here:
a rational projected vector alone would not be a specified Galois
element.

\paragraph{5. The complete middle-weight calculation.}
Write $b_w$ for the standard Lyndon bracket of the word $w$ in the
ordered alphabet $x<y$. The degree-seven bidegree$(4,3)$ words are
\[
 (xxxxyyy,\ xxxyxyy,\ xxxyyxy,\ xxyxxyy,\ xxyxyxy),
\]
and the bidegree$(3,4)$ words are
\[
 (xxxyyyy,\ xxyxyyy,\ xxyyxyy,\ xxyyyxy,\ xyxyxyy).
\]
In these two ordered bases respectively, define
\[
\begin{array}{c|rr}
 &D_1&D_2\\\hline
 D(x)&(1,-5,-6,1,-2)&(-3,-9,-6,-3,-2)\\
 D(y)&(-1,4,6,8,0)&(3,12,6,0,8).
\end{array}
\]
Both derivations are special by direct expansion. They span the
middle-weight special space: the two source spaces have dimension5
each, the target $L_8(4,4)$ has dimension8, and
\[
 (a,b)\longmapsto[a,y]+[x,b]
\]
is surjective. The latter follows from
$L_8=[x,L_7]+[y,L_7]$, itself the Jacobi-identity consequence of
generation by $x,y$, after taking bidegree$(4,4)$.

The degree-six middle Lyndon words are
$(xxxyyy,xxyxyy,xxyyxy)$. Appending their three brackets with $x$
to the columns $D_1(x),D_2(x)$ gives the matrix
\[
 \begin{pmatrix}
 1&-3&1&0&0\\-5&-9&0&1&0\\-6&-6&0&0&1\\
 1&-3&0&0&0\\-2&-2&0&0&0
 \end{pmatrix},\qquad\det=-8.
\]
For $y$ the corresponding matrix is
\[
 \begin{pmatrix}
 -1&3&-1&0&0\\4&12&-1&-1&0\\6&6&0&0&-1\\
 8&0&1&0&-1\\0&8&0&-1&1
 \end{pmatrix},\qquad\det=-8.
\]
Therefore every nonzero middle-weight special $D$ satisfies
\[
 D(x)\notin[x,L_6],\qquad D(y)\notin[y,L_6].
\]
The determinants are nonzero over $\mathbf Q_\ell$ also at
$\ell=2$. No assertion of invertibility over $\mathbf Z_2$ is used.

\paragraph{6. The actual normalizer conjugacy equation through degree seven.}
Suppose $d=1$ and $\log u=ax+$ higher terms, $a\ne0$; the other
line is identical using the second matrix. Lift the $\sigma$ from
Step4 to $n\in N$. Since $\sigma$ is the identity on $V$, its
normalizer exponent is1. Its actual action centralizes $u$.
The canonical boundary-fixing representative differs by an inner
automorphism, so in the rational completed free Lie algebra
\[
 \widetilde\sigma(\log u)
 =\exp(\operatorname{ad}W)(\log u).
\]
Comparison in degree2 gives $W_1=bx$. Replacing the conjugating
element $\exp W$ by
$\exp W\exp(-(b/a)\log u)$ preserves its conjugation action and
removes its degree-one term. Comparison in degrees3,4,5,6 then
successively gives $W_2=W_3=W_4=W_5=0$, using the exact centralizer
$C_L(x)=\mathbf Q_\ell x$. This holds without primitivity of the
original abelian image.

In degree7 the resulting equality is
\[
 aD(x)=[W_6,ax],
\]
contradicting Step5. The second matrix treats $\log u=ay+$ higher
terms. This excludes degree one in every remaining split-CM case.

\paragraph{7. Degree two concludes in the actual pro-$\ell$ group.}
Only $d=2$ remains. The one-dimensional space $L_2(V)$ is
$\det V=\mathbf Q_\ell(1)$. Canceling the nonzero first term of
$\log u$ identifies the actual normalizer character with
$\chi_\ell\circ p$. Nakamura1994,Theorem2.1.1,p.94, now applies
to the class of finite $\ell$-groups and places $A$ in actual cusp
inertia. The finite extension used to define endomorphisms leaves
this geometric cusp subgroup unchanged.

\paragraph{Scope.}
The theorem extends the family of punctured elliptic targets available
in F.3.16 and uses an opened original degree-six source. It does not
construct a detecting elliptic morphism for an arbitrary finite
cover in A.22, and does not treat the higher-dimensional CM quotient
constructed in F.3.19. No conclusion about the full profinite
common-kernel subgroup is drawn.
